\chapter{Conclusiones}

El an\'alisis delimita un sistema de pr\'estamos de bienes de la \escuela{}, no un sistema general de equipos universitarios ni de laboratorios. Su inventario prestable comprende libros, mesa de ping-pong, visor 3D, parlantes y carritos para rob\'otica habilitados por la Escuela. Las computadoras y cualquier recurso adscrito a laboratorios quedan expresamente excluidos.

Los requisitos cubren los puntos de los apuntes de clase: usuarios y perfiles, tiempos y pol\'iticas, reservas, registros, inventario, estados del pr\'estamo, historial, garant\'ias y sanciones. Se diferencia el estado f\'isico de una unidad de la disponibilidad de un intervalo y del estado de un pr\'estamo; de este modo, una reserva futura no equivale a una entrega ni un vencimiento libera el bien.

El diagrama UML distingue tipos de bienes, fichas de cat\'alogo y unidades f\'isicas, y organiza el negocio en contextos delimitados. Las reglas propuestas previenen dobles entregas, reutilizaci\'on de reservas, p\'erdida de trazabilidad y cambios retroactivos de condiciones. La extensi\'on UML en paquetes y clases ubica estos conceptos en el dominio y muestra c\'omo los casos de uso, los controladores y los adaptadores colaboran sin mezclar responsabilidades.

Antes del dise\~no e implementaci\'on definitiva se deben validar con la Escuela los plazos, cupos, modalidades de uso, garant\'ias, sanciones y requisitos de despliegue.
